\clearpage\begingroup\NoSourceHead\SourceLeaf{237}
\begin{center}
[II.]\par\medskip
[WEITERES ZUR BAHNBESTIMMUNG.]\par\medskip
\rule{0.92\linewidth}{0.35pt}\par
{\small [Aus Scheda Ag, Cereri Ferdinandeae Sacrum; November 1801, S. 5—17.]}\par
\rule{0.92\linewidth}{0.35pt}\par\medskip
[1.]\par\medskip
\textls[100]{Vorschriften,}\\um aus
\end{center}
\noindent{\small [S. 5]}\par
\begin{center}
\begin{tabular}{@{}l@{\quad}l@{\quad}l@{}}
geocent[rischer] Länge eines Weltkörpers & . . . . . . . . & $\alpha$\\
\hspace{3em}Breite & . . . . . . . . . . . . . . . . . . . & $\beta$\\
Neigung seiner Bahn & . . . . . . . . . . . . . . . . . . . . & $i$\\
Ort des aufsteigenden Knoten & . . . . . . . . . . . . . . . & $\AstNode$\\
Länge der Sonne & . . . . . . . . . . . . . . . . . . . . . . & $V$\\
Abstand der Sonne von der Erde & . . . . . . . . . . . . . & $R$
\end{tabular}\par\medskip
\textls[100]{abzuleiten:}\par\medskip
\begin{tabular}{@{}l@{\quad}l@{\quad}l@{}}
heliocentrische Länge des Weltkörpers & . . . . . . . . . . & $v$\\
Abstand desselben von der Sonne & . . . . . . . . . . . . & $r$\\
Curtirter Abstand von der Erde & . . . . . . . . . . . . . & $\Delta$.
\end{tabular}
\end{center}
\[
\mathrm{I.}\quad\operatorname{cotang}(v-\AstNode)=\frac{\cos i\cos(V-\AstNode)\operatorname{tang}\beta-\sin i\sin(\alpha-V)}{\sin(V-\AstNode)\operatorname{tang}\beta},
\]
\[
\mathrm{II.}\quad\Delta=R\frac{\operatorname{tg}i\sin(V-\AstNode)}{\operatorname{tg}i\sin(\alpha-\AstNode)-\operatorname{tg}\beta},
\]
\[
\mathrm{III.}\quad r=\Delta\frac{\operatorname{tang}\beta}{\sin(v-\AstNode)\sin i},
\]
\noindent [Zur Kontrolle:]
\[
\operatorname{tang}(v-\AstNode)=\frac1{\cos i}\mathbin{.}\frac{\Delta\sin(\alpha-\AstNode)-R\sin(V-\AstNode)}{\Delta\cos(\alpha-\AstNode)-R\cos(V-\AstNode)},
\]
\[
r=\surd\bigl(\Delta\Delta\sec\beta^2+RR-2R\Delta\cos(\alpha-V)\bigr).
\]
\clearpage
\HistoricalHead{233}{R}{WEITERES ZUR BAHNBESTIMMUNG.}\SourceLeaf{238}\HistoricalFoot{XI1.}{30}
\begin{center}
\textls[100]{Beispiele}[\textsuperscript{*})].\par\medskip
I. [erste Beobachtung.]
\end{center}
\noindent{\small [S. 6]}\par
\begingroup\renewcommand{\thefootnote}{}\footnotetext{[\textsuperscript{*}) Diese Beispiele beziehen sich auf \textsc{Piazzis} Beobachtungen der Ceres, vergl. S. 239.]}\endgroup
\[
\begin{array}{r@{\;}c@{\;}l@{\qquad}r@{\;}c@{\;}l}
V&=&282^0\;2'\;28''60&\log R&=&9{,}992\,6317\\
\alpha&=&53\;19\;42''35&\log\operatorname{tg}\beta&=&8{,}725\,0840n\\
\AstNode&=&81\;27\;4''62&\log\operatorname{tg}i&=&9{,}262\,1790
\end{array}
\]
\[
\begin{aligned}
V-\AstNode&=200^0\;35'\;23''98\\
\alpha-\AstNode&=331\;52\;37''73\\
\alpha-V&=131\;17\;13''75
\end{aligned}
\]
\[
\begin{array}{r@{\;}c@{\;}l}
\log\cos i&=&9{,}992\,8560\\
\log\operatorname{cot}(V-\AstNode)&=&0{,}425\,1865\\\cline{2-3}
&&0{,}418\,0426
\end{array}
\]
\[
\begin{array}{r@{\;}c@{\;}l}
\log\sin i&=&9{,}255\,0351\\
\log\sin(\alpha-V)&=&9{,}875\,8782\\
\mathrm{Compl.}\,\mathrm{l[og]}\sin(V-\AstNode)&=&0{,}453\,8546n\\
\mathrm{Compl.}\,\mathrm{l[og]}\operatorname{tang}\beta&=&1{,}274\,9160n\\\cline{2-3}
&&0{,}859\,6839
\end{array}
\]
\[
\begin{array}{r@{\quad}l}
{}[\cos i\mathbin{.}\operatorname{cotang}(V-\AstNode)=]&2{,}618\,4398\\[1ex]
\left[\dfrac{\sin i\mathbin{.}\sin(\alpha-V)}{\sin(V-\AstNode)\mathbin{.}\operatorname{tang}\beta}=\right]&7{,}239\,0833\text{[\textsuperscript{**})]}\\\cline{2-2}
{}[\operatorname{cotang}(v-\AstNode)=-]&4{,}620\,6435\\
{}[\log=]&0{,}664\,7025[n]
\end{array}
\]
\begingroup\renewcommand{\thefootnote}{}\footnotetext{[\textsuperscript{**}) Der richtige Wert ist bei Anwendung 7-stelliger Tafeln $7{,}239\,0883$.]}\endgroup
\[
\begin{array}{r@{\quad}l}
v-\AstNode=&-12^0\;12'\;41''85\\
{}[\AstNode=]&81\;27\;4\;62\\\cline{2-2}
v=&69\;14\;22\;77
\end{array}
\]
\begingroup\small\setlength{\arraycolsep}{1pt}
\[
\begin{array}{r@{\;}l@{\;}r@{\;}l}
\sin(\alpha-\AstNode)=&\text{. . . . . . . . . . . . . . .}\;-0{,}471\,36368&&\\[1ex]
\log\dfrac{\operatorname{tg}\beta}{\operatorname{tg}i}=&9{,}462\,9050n\;\text{. . .}\;\underline{-0{,}290\,33867\text{[\textsuperscript{***})]}}&\log\sin(V-\AstNode)=&9{,}546\,1454n\\
\text{[Differenz }=]&-0{,}181\,02501\;\text{. . . . . . . . . . . . .}&\log=&9{,}257\,7386n\\\cline{4-4}
&&&0{,}288\,4068\\[1ex]
&&[\log R=]&9{,}992\,6317\\\cline{4-4}
&&\log\Delta=&0{,}281\,0385
\end{array}
\]
\endgroup
\begingroup\renewcommand{\thefootnote}{}\footnotetext{[\textsuperscript{***}) Der genaue Wert ist bei Anwendung 7-stelliger Tafeln $0{,}290\,33873$.]}\endgroup
\clearpage
\HistoricalHead{234}{L}{THEORETISCHE ASTRONOMIE. NACHLASS.}\SourceLeaf{239}
\[
\begin{array}{r@{\;}c@{\;}l}
\log\Delta&=&0{,}281\,0385\\
\log\operatorname{tg}\beta&=&8{,}725\,0840n\\
\mathrm{C.}\log\sin i&=&0{,}744\,9649\\
\mathrm{C.}\log\sin(v-\AstNode)&=&0{,}674\,6422n\\\cline{2-3}
\log r&=&0{,}425\,7296
\end{array}
\]
\begin{center}II. [zweite Beobachtung].\end{center}
\noindent{\small [S. 7]}\par
\[
\begin{array}{r@{\;}c@{\;}l@{\qquad}r@{\;}c@{\;}l}
V&=&322^0\;35'\;5''89&\log R&=&9{,}994\,5823\\
\alpha&=&56\;26\;43''70&\log\operatorname{tg}\beta&=&8{,}020\,1510n\\
\AstNode&=&81\;27\;4''62&\log\operatorname{tg}i&=&9{,}262\,1790
\end{array}
\]
\noindent [Die Rechnung ergab]
\[
\begin{array}{r@{\;}l}
{}[v=]&78^0\;26'\;32''49\\
\log\Delta=&0{,}374\,1524\\
\log r=&0{,}419\,1792.
\end{array}
\]
\begin{center}III. [dritte Beobachtung].\end{center}
\noindent{\small [S. 8]}\par
\[
\begin{array}{r@{\;}c@{\;}l@{\qquad}r@{\;}c@{\;}l}
V&=&302^0\;21'\;29''24&\log R&=&9{,}993\,1886\\
\alpha&=&53\;39\;2''62&\log\operatorname{tg}\beta&=&8{,}474\,2601\\
\AstNode&=&81\;27\;4''62&\log\operatorname{tg}i&=&9{,}262\,1790.
\end{array}
\]
\noindent [Die Rechnung ergab]
\[
\begin{array}{r@{\;}l}
{}[v=]&73^0\;48'\;36''17\text{[\textsuperscript{*})]}\\
\log\Delta=&0{,}327\,2519\\
{}[\log r=]&0{,}422\,7241\text{[\textsuperscript{**})]}.
\end{array}
\]
\begingroup\renewcommand{\thefootnote}{}
\footnotetext{[\textsuperscript{*}) Später verbessert in] $35''77$.}
\footnotetext{[\textsuperscript{**}) Die letzten drei Dezimalen später verbessert in] 177.}
\endgroup
\clearpage\endgroup
